\section{USAMO 2018}
        \begin{enumerate}
        	\item (\textit{USAMO 2018}) Note: For any geometry problem whose statement begins with an asterisk ($*$), the first page of the solution must be a large, in-scale, clearly labeled diagram. Failure to meet this requirement will result in an automatic 1-point deduction.


 Let $a,b,c$ be positive real numbers such that $a+b+c=4\sqrt[3]{abc}$. Prove that 
 \[2(ab+bc+ca)+4\min(a^2,b^2,c^2)\ge a^2+b^2+c^2.\]


 



\textbf{Solution 1}

\href{https://wiki-images.artofproblemsolving.com//6/69/IMG\_8946.jpg}{https://wiki-images.artofproblemsolving.com//6/69/IMG\_8946.jpg}

 -srisainandan6

 



\textbf{Solution 2}

Similarly to Solution 2, we will prove homogeneity but we will use that to solve the problem differently. Let $f(a,b,c)=a+b+c-4\sqrt[3]{abc}$. Note that $f(a,b,c)=f(ka,kb,kc)$, thus proving homogeneity. 

 WLOG, we can scale down all variables such that the lowest one is $1$. WLOG, let this be $a=1$. We now have $1+b+c=4\sqrt[3]{bc}$, and we want to prove $2bc+2b+2c+4\ge 1+b^2+c^2.$ Adding $2bc$ to both sides and subtracting $2b+2c$ gives us $4bc+4\ge 1+ (b+c)(b+c-2)$, or $4bc+3\ge (b+c)(b+c-2)$. Let $\sqrt[3]{bc}=x$. Now, we have 
 \[4x^3+3 \ge (4x-1)(4x-3)\] 
 \[4x^3 - 16x^2 + 16x \ge 0\] 
 \[4x^2 - 16 + 16 \ge 0\] 
 \[4(x-2)^2 \ge 0\] By the trivial inequality, this is always true. Since all these steps are reversible, the proof is complete. ~SigmaPiE

 



\textbf{Solution 3}

WLOG, let $a \le b$ and $a \le c$ and add $2(ab+bc+ca)$ to both sides to make the left side a square.
 \[4(ab+bc+ca)+4a^2 \ge (a+b+c)^2\] 
 \[4bc+4a(a+b+c) \ge (a+b+c)^2\] Now we perform the substitution for $a+b+c$. 
 \[4bc+4a(4(abc)^{1/3}) \ge 16(abc)^{2/3}\] Multiply both sides by $a$ because we want all the terms to have a factor of $(abc)^{1/3}$. 
 \[4abc+16a^2(abc)^{1/3} \ge 16a(abc)^{2/3}\] Divide both sides by $4(abc)^{1/3}$ . 
 \[(abc)^{2/3}+4a^2 \ge 4a(abc)^{1/3}\] 
 \[((abc)^{1/3}-2a)^2 \ge 0\] This is true by the trivial inequality. Lastly, all the steps are reversible so the given inequality has been proved.

 -Themathcanadian

 


	\item (\textit{USAMO 2018}) Note: For any geometry problem whose statement begins with an asterisk ($*$), the first page of the solution must be a large, in-scale, clearly labeled diagram. Failure to meet this requirement will result in an automatic 1-point deduction.


 Find all functions $f:(0,\infty) \rightarrow (0,\infty)$ such that


 
 \[f\left(x+\frac{1}{y}\right)+f\left(y+\frac{1}{z}\right) + f\left(z+\frac{1}{x}\right) = 1\]
for all $x,y,z >0$ with $xyz =1.$


 



\textbf{Solution}

Obviously, the output of $f$ lies in the interval $(0,1)$. Define $g:(0,1)\to(0,1)$ as $g(x)=f\left(\frac1x-1\right)$. Then for any $a,b,c\in(0,1)$ such that $a+b+c=1$, we have $g(a)=f\left(\frac1a-1\right)=f\left(\frac{1-a}a\right)=f\left(\frac{b+c}a\right)$. We can transform $g(b)$ and $g(c)$ similarly:

 
 \[g(a)+g(b)+g(c)=f\left(\frac ca+\frac ba\right)+f\left(\frac ab+\frac cb\right)+f\left(\frac bc+\frac ac\right)\]

 Let $x=\frac ca$, $y=\frac ab$, $z=\frac bc$. We can see that the above expression is equal to $1$. That is, for any $a,b,c\in(0,1)$ such that $a+b+c=1$, $g(a)+g(b)+g(c)=1$.

 (To motivate this, one can start by writing $x=\frac ab$, $y=\frac bc$, $z=\frac ca$, and normalizing such that $a+b+c=1$.)

 
For convenience, we define $h:\left(-\frac13,\frac23\right)\to\left(-\frac13,\frac23\right)$ as $h(x)=g\left(x+\frac13\right)-\frac13$, so that for any $a,b,c\in\left(-\frac13,\frac23\right)$ such that $a+b+c=0$, we have 

 
 \[h(a)+h(b)+h(c)=g\left(a+\frac13\right)-\frac13+g\left(b+\frac13\right)-\frac13+g\left(c+\frac13\right)-\frac13=1-1=0.\]

 Obviously, $h(0)=0$. If $|a|<\frac13$, then $h(a)+h(-a)+h(0)=0$ and thus $h(-a)=-h(a)$. Furthermore, if $a,b$ are in the domain and $|a+b|<\frac13$, then $h(a)+h(b)+h(-(a+b))=0$ and thus $h(a+b)=h(a)+h(b)$. 

 
At this point, we should realize that $h$ should be of the form $h(x)=kx$. We first prove this for some rational numbers. If $n$ is a positive integer and $x$ is a real number such that $|nx|<\frac13$, then we can repeatedly apply $h(a+b)=h(a)+h(b)$ to obtain $h(nx)=nh(x)$. Let $k=6h\left(\frac16\right)$, then for any rational number $r=\frac pq\in\left(0,\frac13\right)$ where $p,q$ are positive integers, we have $h(r)=6p*h\left(\frac1{6q}\right)=\frac{6p}q*h\left(\frac16\right)=kr$.

 Next, we prove it for all real numbers in the interval $\left(0,\frac13\right)$. For the sake of contradiction, assume that there is some $x\in\left(0,\frac13\right)$ such that $h(x)\ne kx$. Let $E=h(x)-kx$, then obviously $0<|E|<1$. The idea is to "amplify" this error until it becomes so big as to contradict the bounds on the output of $h$. Let $N=\left\lceil\frac1{|E|}\right\rceil$, so that $N\ge2$ and $|NE|\ge1$. Pick any rational $r\in\left(\frac{N-1}Nx,x\right)$, so that 
 \[0<x-r<N(x-r)<x<\frac13.\]All numbers and sums are safely inside the bounds of $\left(-\frac13,\frac13\right)$. Thus
 \[h(N(x-r))=Nh(x-r)=N(h(x)+h(-r))=N(h(x)-h(r))=kN(x-r)+NE,\]but picking any rational number $s\in\left(N(x-r),\frac13\right)$ gives us $|kN(x-r)|<|ks|$, and since $ks=h(s)\in\left(-\frac13,\frac23\right)$, we have $kN(x-r)\in\left(-\frac13,\frac23\right)$ as well, but since $NE\ge1$, this means that $h(N(x-r))=kN(x-r)+NE\notin\left(-\frac13,\frac23\right)$, giving us the desired contradiction.

 
We now know that $h(x)=kx$ for all $0<x<\frac13$. Since $h(-x)=-h(x)$ for $|x|<\frac13$, we obtain $h(x)=kx$ for all $|x|<\frac13$. For $x\in\left(\frac13,\frac23\right)$, we have $h(x)+h\left(-\frac x2\right)+h\left(-\frac x2\right)=0$, and thus $h(x)=kx$ as well. So $h(x)=kx$ for all $x$ in the domain. Since $h(x)$ is bounded by $-\frac13$ and $\frac23$, we have $-\frac12\le k\le1$. It remains to work backwards to find $f(x)$.

 
 \begin{align*} h(x) &= kx \\ g(x) &= kx+\frac{1-k}3 \\ f(x) &= \frac k{1+x}+\frac{1-k}3\quad\left(-\frac12\le k\le1\right). \end{align*}
- wzs26843545602

 


	\item (\textit{USAMO 2018}) Note: For any geometry problem whose statement begins with an asterisk ($*$), the first page of the solution must be a large, in-scale, clearly labeled diagram. Failure to meet this requirement will result in an automatic 1-point deduction.


 For a given integer $n\ge 2,$ let $\{a_1,a_2,\ldots,a_m\}$ be the set of positive integers less than $n$ that are relatively prime to $n.$ Prove that if every prime that divides $m$ also divides $n,$ then $a_1^k+a_2^k + \dots + a_m^k$ is divisible by $m$ for every positive integer $k.$


 



\textbf{Solution (NO FORMATTING)}

\href{https://maa.org/sites/default/files/pdf/AMC/usamo/2018/2018USAMO.pdf}{https://maa.org/sites/default/files/pdf/AMC/usamo/2018/2018USAMO.pdf}The integer m in the statement of the problem is $\phi$(n), where $\phi$ is the Euler totient function.Throughout our proof we write ps|| m, if s is the greatest power of p that divides m.We begin with the following lemma:Lemma 1. If p is a prime and ps divides n for some positive integer s, then 1k + 2k + · · · + nkisdivisible by ps-1for any integer k $\geq$ 1.2018 USAMO – Solutions 4Proof. Let {a1, a2, . . . , am} be a complete reduced residue set modulo ps and m = ps-1(p - 1).First we prove by induction on s that for any positive integer k, ak1 + ak2 + · · · + akm is divisible byps-1. The base case s = 1 is true. Suppose the statement holds for some value of s. Consider thestatement for s + 1. Note that{a1, . . . , am, ps + a1, . . . , ps + am, . . . , ps(p - 1) + a1, . . . , ps(p - 1) + am}is a complete reduced residue set modulo ps+1. Therefore, the desired sum of k-th powers is equaltoak1 + · · · + akm + · · · + (ps(p - 1) + a1)k + · · · + (ps(p - 1) + am)k \equiv p(ak1 + · · · + akm) \equiv 0 (mod ps),where we have used the induction hypothesis for the second congruence. This gives the inductionstep.Now we are ready to prove the lemma. Because numbers from 1 to n can be split into blocks ofconsecutive numbers of length ps, it is enough to show that 1k + 2k +· · · + (ps)kis divisible by ps-1for any positive integer k. We use induction on s. The statement is true for s = 1. Assume thestatement is true for s - 1. The sum1k + 2k + · · · + (ps)k = ak1 + ak2 + · · · + akm + pk1k + 2k + · · · + (ps-1)kis divisible by ps-1, because ps-1| ak1 + · · · + akm and by the induction hypothesis ps-2| 1k + · · · +(ps-1)k.Now we proceed to prove a second lemma, from which the statement of the problem will immediatelyfollow:Lemma 2. Suppose p is a prime dividing n. Let {a1, . . . , am} be a complete reduced residue setmod n, and define s by ps|| m. Then ps divides ak1 + · · · + akm for any integer k $\geq$ 1.Proof. We fix p, and use induction on the number of prime factors of n (counted by multiplicity)that are different from p. If there are no prime factors other than p, then n = ps+1, m = ps(p - 1),and we proved in Lemma 1 that ak1 +· · ·+akm is divisible by ps. Now suppose the statement is truefor n. We show that it is true for nq, where q is a prime not equal to p.Case 1. q divides n. We have ps|| $\phi$(n) and ps|| $\phi$(nq), because $\phi$(nq) = q$\phi$(n). If {a1, a2, . . . , am}is a complete reduced residue set modulo n, then{a1, . . . , am, n + a1, . . . , n + am, . . . , n(q - 1) + a1, . . . , n(q - 1) + am}is a complete reduced residue set modulo nq. The new sum of k-th powers is equal toak1 + · · · + akm + · · · + (n(q - 1) + a1)k + · · · + (n(q - 1) + am)k = mnk1k + · · · + (q - 1)k+2018 USAMO – Solutions 5k1nk-11k-1 + · · · + (q - 1)k-1(a1 + · · · + am) + · · · + q(ak1 + · · · + akm).This sum is divisible by ps because ps|| m and ps| aj1 + aj2 + · · · + ajm for any positive integer j.Case 2. q doesn’t divide n. Suppose pb|| q - 1, where b $\geq$ 0. Note that $\phi$(nq) = $\phi$(n)(q - 1), sops|| $\phi$(n) and ps+b|| $\phi$(nq). Let {a1, . . . , am} be a complete reduced residue set modulo n. Thecomplete reduced residue set modulo nq consists of the mq numbers{a1, . . . , am, n + a1, . . . , n + am, . . . , n(q - 1) + a1, . . . , n(q - 1) + am}with the m elements {qa1, qa2, . . . , qam} removed.The new sum of k-th powers is equal toak1 + · · · + akm + · · · + (n(q - 1) + a1)k + · · · + (n(q - 1) + am)k - qk(ak1 + · · · + akm) =mnk1k + · · · + (q - 1)k+k1nk-11k-1 + · · · + (q - 1)k-1(a1 + · · · + am) + · · ·· · · +kk - 1n (1 + · · · + (q - 1)) (ak-11 + · · · + ak-1m ) + q(ak1 + · · · + akm) - qk(ak1 + · · · + akm).Each termkjnk-j1k-j + · · · + (q - 1)k-j(aj1 + · · · + ajm),for 0 $\leq$ j $\leq$ k - 1, is divisible by ps+b because p | nk-j, ps| aj1 + · · · + ajm, and pb-1| 1k-j + · · · +(q - 1)k-j by Lemma 1.Also (qk - q)(ak1 + · · · + akm) is divisible by ps+b because pb| q - 1 | qk - q and ps| ak1 + · · · + akm.Thus ps+b divides our sum and our proof is complete.Remark. In fact, one can also show the converse statement: if {a1, a2, . . . , am} is as defined inthe problem and ak1 + ak2 + · · · + akm is divisible by m for every positive integer k, then every primethat divides m also divides n.

 


	\item (\textit{USAMO 2018}) Note: For any geometry problem whose statement begins with an asterisk ($*$), the first page of the solution must be a large, in-scale, clearly labeled diagram. Failure to meet this requirement will result in an automatic 1-point deduction.


 Let $p$ be a prime, and let $a_1, \dots, a_p$ be integers. Show that there exists an integer $k$ such that the numbers 
 \[a_1 + k, a_2 + 2k, \dots, a_p + pk\]produce at least $\frac{1}{2} p$ distinct remainders upon division by $p$.


 



\textbf{Solution}

$\textbf{Lemma: }$ For fixed $i\neq j,$ where $i, j\in\{1, 2, ..., p\},$ the statement $a_i + ik\equiv a_j + jk\text{ (mod } p\text{)}$ holds for exactly one $k\in {1, 2, ..., p}.$

 $\textbf{Proof: }$ Notice that the left side minus the right side is congruent to $(a_i - a_j) + (i - j)k$ modulo $p.$ For this difference to equal $0,$ there is a unique solution for $k$ modulo $p$ given by $k\equiv (a_j - a_i)(i - j)^{-1}\text{ (mod } p\text{)},$ where we have used the fact that every nonzero residue modulo $p$ has a unique multiplicative inverse. Therefore, there is exactly one $k\in {1, 2, ..., p}$ that satisfies $a_i + ik\equiv a_j + jk\text{ (mod } p\text{)}$ for any fixed $i\neq j. \textbf{ End Lemma}$

 Suppose that you have $p$ graphs $G_1, G_2, ..., G_p,$ and graph $G_k$ consists of the vertices $(i, k)$ for all $1\le i\le p.$ Within any graph $G_k,$ vertices $(i_1, k)$ and $(i_2, k)$ are connected by an edge if and only if $a_{i_1} + i_1k\equiv a_{i_2} + i_2k\text{ (mod } p\text{)}.$ Notice that the number of disconnected components of any graph $G_k$ equals the number of distinct remainders when divided by $p$ given by the numbers $a_1 + k, a_2 + 2k, ..., a_p + pk.$

 These $p$ graphs together have exactly one edge for every unordered pair of elements of $\{1, 2, ..., p\},$ so they have a total of exactly $\frac{p(p-1)}{2}$ edges. Therefore, there exists at least one graph $G_k$ that has strictly fewer than $\frac{p}{2}$ edges, meaning that it has more than $\frac{p}{2}$ disconnected components. Therefore, the collection of numbers $\{a_i + ik: 1\le i\le p\}$ for this particular value of $k$ has at least $\frac{p}{2}$ distinct remainders modulo $p.$ This completes the proof.

 $\textbf{Note: This is the same problem as 2018 USAJMO Problem 5.}$

 



\textbf{Solution 2}

We consider the Lemma and it's proof as in the above solution.

 For every $k$, let $x_k$ be the number of distinct residues that $a_i + ik$ takes on. Further for each residue $r_i$, let $n_i$ be the number of times it is achieved by an $a_j + jk$ for $i \in [1, x_k]$. Note that $\sum_{i = 1}^{x_k} n_i = p.$

 The number of pairs $(a_i, a_j)$ s.t $(a_i + ik) - (a_j + jk) \equiv 0$ for a given $k$ is,
 \begin{align*} \sum_{i = 1}^{x_k} \binom{n_i}{2} &\ge x_k \binom{(\sum_{i = 1}^{x_k} n_i)/x_k}{2} \\ &= \frac{1}{2}p(\frac{p}{x_k} - 1). \end{align*}
The first equality is given by a well known theorem, which can be proven by C-S or Jensen's.

 Summing over all $k$, the number of times a pair $(a_i, a_j)$ has $(a_i + ik) - (a_j + jk) \equiv 0$ is,
 \begin{align*} \sum_{k = 1}^{p} \frac{1}{2}p(\frac{p}{x_k} - 1). \end{align*}
Alternatively, every pair $(a_i, a_j)$ has a unique $k$ s.t $(a_i + ik) - (a_j + jk) \equiv 0$ by the Lemma so we double-count as,
 \begin{align*} \binom{p}{2} &\ge \sum_{k = 1}^{p} \frac{1}{2}p(\frac{p}{x_k} - 1). \\ 2p - 1 &\ge \sum_{k = 1}^{p} \frac{p}{x_k}. \end{align*}
By AM-HM inequality,
 \begin{align*} \sum_{k = 1}^{p} \frac{p}{x_k} &\ge \frac{p^2}{\sum_{k = 1}^{p} \frac{x_k}{p}} \\ 2p-1 &\ge \frac{p^3}{\sum_{k = 1}^{p} x_k} \\ \sum_{k = 1}^{p}x_k &\ge \frac{p^3}{2p-1} \ge \frac{p^2}{2}. \end{align*}
So it is clear that $\exists k$ s.t $x_k \ge \frac{p}{2}$.

 


	\item (\textit{USAMO 2018}) Note: For any geometry problem whose statement begins with an asterisk ($*$), the first page of the solution must be a large, in-scale, clearly labeled diagram. Failure to meet this requirement will result in an automatic 1-point deduction.


 In convex cyclic quadrilateral $ABCD,$ we know that lines $AC$ and $BD$ intersect at $E,$ lines $AB$ and $CD$ intersect at $F,$ and lines $BC$ and $DA$ intersect at $G.$ Suppose that the circumcircle of $\triangle ABE$ intersects line $CB$ at $B$ and $P$, and the circumcircle of $\triangle ADE$ intersects line $CD$ at $D$ and $Q$, where $C,B,P,G$ and $C,Q,D,F$ are collinear in that order. Prove that if lines $FP$ and $GQ$ intersect at $M$, then $\angle MAC = 90^{\circ}.$


 



\textbf{Solution}


                    \begin{figure}[H]
                        \centering
                        \includegraphics[width=0.5\textwidth]{Solution/USAMO_I_2018_5_Solution_Figure_1.png}
                    \end{figure}
                    

 
 \begin{align*} &\mathrel{\phantom{=}}\angle DEQ+\angle AED+\angle AEP\\ &=\angle DAQ+\angle AQD+\angle AEP\\ &=180-\angle ADC+\angle AEP\\ &=180-\angle ADC+\angle ABP\\ &=\angle ABP+\angle ABC\\ &=180 \end{align*}
so $P,E,Q$ are collinear. Furthermore, note that $DQBP$ is cyclic because:
 \[\angle EDQ = \angle BAE = BPE.\]Notice that since $A$ is the intersection of $(EDQ)$ and $(BPE)$, it is the Miquel point of $DQBP$. 

 Now define $X$ as the intersection of $BQ$ and $DP$. From Pappus's theorem on $BFPDGQ$ that $A,M,X$ are collinear. It’s a well known property of Miquel points that $\angle EAX = 90$, so it follows that $MA \perp AE$, as desired. $\blacksquare$~AopsUser101

 


	\item (\textit{USAMO 2018}) Note: For any geometry problem whose statement begins with an asterisk ($*$), the first page of the solution must be a large, in-scale, clearly labeled diagram. Failure to meet this requirement will result in an automatic 1-point deduction.


 Let $a_n$ be the number of permutations $(x_1, x_2, \dots, x_n)$ of the numbers $(1,2,\dots, n)$ such that the $n$ ratios $\frac{x_k}{k}$ for $1\le k\le n$ are all distinct. Prove that $a_n$ is odd for all $n\ge 1.$


 



\textbf{Solution}

Write out the mapping of each $k$ to $x_k$ as follows:
 \[1\ \ 2\ \ 3\ \ \dots \ \ n\]
 \[x_1\ x_2\ x_3\ \dots \ x_n\]

 Now, consider what happens when the two rows are swapped (and the top-bottom pairs are reordered so that the top reads (1,2,3,...,n)).  This will result in a new permutation if and only if $x_k = j$ does not imply $x_j = k$ for all $k,j$ in $(1,2,3,\dots,n)$.  I will denote this new permutation as $(y_1,y_2,y_3,\dots,y_n)$.  Notice that $(y_1,y_2,y_3,\dots,y_n)$ is a valid permutation if and only if $(x_1,x_2,x_3,\dots,x_n)$ is valid, because each fraction $\frac{y_k}{k}$ is the reciprocal of $\frac{x_k}{k}$.  This means that we need only consider the parity of number of cases in which $(y_1,y_2,y_3,\dots, y_n) = (x_1,x_2,x_3,\dots, x_n)$, as there will be an even number of cases where $(y_1,y_2,y_3,\dots, y_n) \ne (x_1,x_2,x_3,\dots, x_n)$.  Let the number of valid permutations where $(y_1,y_2,y_3,\dots, y_n) = (x_1,x_2,x_3,\dots, x_n)$ be $b_n$.

 Notice that the only permutations that have the property $x_{x_k}=k$ (which is an equivalent statement to the one given above) are those that are formed by taking pairs of elements and swapping their positions having the maximum number of pairs possible and having no two pairs both contain the same element.  Some more necessary conditions will be outlined below after we split $n$ into two cases.

 
Case 1:  $n$ is odd.  If $n$ is odd then there must be exactly one $k$ such that $x_k = k$.  This yields $b_n=n\cdot b_{n-1}$* which has the same parity as $b_{n-1}$, so we need only consider the parity of $b_n$ for odd $n$. (*\textbf{Note: This is wrong, as $b_n$ is only the number of valid involutions for $1,2,...,n$, not $1,2,...,n+1$ with one number removed.})

 Case 2:  $n$ is even.  If $n$ is even then can be no $k$ so that $x_k=k$, or else there must be at least two distinct numbers $k$ and $j$ so that $x_k=k$ and $x_j=j$ which violates the given conditions.  Denote a pair of numbers that are swapped to arrive at the final permutation as the pair $(a,b)$.  Then if a permutation yields an invalid arrangement there must be two pairs $(a,b)$ and $(c,d)$ such that $\frac{a^2}{b^2} = \frac{c^2}{d^2}$.  But notice that the two pairs $(a,c)$ and $(b,d)$ will also result in an invalid arrangement. So, there must be an even number of invalid arrangements (\textbf{Note: This proof doesn't work, as no specific pairing of invalid arrangements is constructed. Indeed, when three pairs all with the same fraction are present, this doesn't work.}), meaning the parity we desire is just the number of ways to separate $n$ objects into $\frac n2$ pairs!  However, this is quite simply just $(n-1)(n-3)(n-5)\cdots(3)(1)$, which is clearly the product of odd numbers.  So we conclude that there are an odd number of valid permutations $(x_1,x_2,x_3,\dots,x_n)$.  $\blacksquare$

 


\end{enumerate}